% --- Página de copyright ---
% La portada principal la genera pandoc con \maketitle (metadata
% title/subtitle/author/date). Aquí solo agregamos las páginas que
% pandoc no produce nativamente: copyright + dedicatoria.
\clearpage
\thispagestyle{empty}
\vspace*{2cm}

\noindent\textbf{AgencyDomains}\\
Arquitectura del Mundo Agentivo\\
César Obach-Renner

\vspace{1cm}

\noindent\textbf{Editor:} GegoLabs\\
\textbf{Edición:} Borrador de desarrollo · Julio 2026 (v0.7)\\
\textbf{Licencia:} GNU Free Documentation License v1.3 (propuesta)

\vspace{1cm}

\noindent Este libro se publica bajo GNU Free Documentation License v1.3 (propuesta). El lector puede copiar, distribuir y modificar la obra bajo los términos de la licencia. La sección invariante es el Prefacio.

\vspace{0.8cm}

\noindent\textbf{Cómo citar:}\\
Obach-Renner, César. \emph{AgencyDomains: arquitectura del Mundo Agentivo}. Borrador de desarrollo v0.7. GegoLabs, 2026.

\vspace{0.8cm}

\noindent\textbf{Información de contacto:}\\
Sitio web: agencydomains.org\\
Repositorio: github.com/gegolabs/agencydomains.org\\
Errata y comentarios: github.com/gegolabs/agencydomains.org/issues

\vfill

% --- Dedicatoria ---
\clearpage
\thispagestyle{empty}
\vspace*{\fill}
\begin{center}
\textit{A ti, Alejandra:\\
estos tres libros hablan de un viaje hacia el futuro,\\
pero el viaje que me importa es el que hago contigo.\\
Gracias por sostener, todos estos años,\\
mi búsqueda incesante por crear tecnología.\\[0.8em]
Y también a nuestros hijos, Matías, Victoria y Sebastián,\\
que han crecido al lado de esa búsqueda.}
\end{center}
\vspace*{\fill}
\clearpage
